\documentclass{statement}
\usepackage{tikz}
\usepackage{ulem}
\usepackage{tabularx}
\usepackage{makecell}
\usepackage{tabularray}
\usepackage{color}
\usepackage{xcolor}
\usepackage{hyperref}
\usepackage{minted}

\title{ 示例比赛 }
\date{}

\begin{document}
\begin{titlepage}
\vspace*{-20mm}
\begin{center}
{\fontsize{22}{32}\selectfont \heiti 示例比赛}\\
\vskip 0.6em
\fontsize{15}{22}\selectfont \rmfamily 时间：2026 年 8 月 3 日 09:00 ～ 13:00（4 小时）
\vskip 0.5em
\vskip 0.8em
\end{center}

\begin{center}
\begin{tabularx}{\textwidth}{|l|>{\centering\arraybackslash}X|>{\centering\arraybackslash}X|>{\centering\arraybackslash}X|>{\centering\arraybackslash}X|}
\hline
    题目名称 & 数列之美 & 简单求和 & 棋盘平方 & 回文判断 \\
\hline
    题目类型 & 传统型 & 传统型 & 传统型 & 传统型 \\
\hline
    目录 & \texttt{seq} & \texttt{demo2} & \texttt{demo3} & \texttt{demo4} \\
\hline
    可执行文件名 & \texttt{seq} & \texttt{demo2} & \texttt{demo3} & \texttt{demo4} \\
\hline
    输入文件名 & \texttt{seq.in} & 标准输入 & 标准输入 & 标准输入 \\
\hline
    输出文件名 & \texttt{seq.out} & 标准输出 & 标准输出 & 标准输出 \\
\hline
    每个测试点时限 & 1.0 秒 & 1.0 秒 & 2.0 秒 & 0.5 秒 \\
\hline
    内存限制 & 512 MiB & 256 MiB & 512 MiB & 128 MiB \\
\hline
    测试点数目 & 1 & 1 & 1 & 1 \\
\hline
\end{tabularx}
\end{center}

{\noindent\hspace*{1.5em}\rmfamily 提交源程序文件名}\par
\begin{center}
\begin{tabularx}{\textwidth}{|l|>{\centering\arraybackslash}X|>{\centering\arraybackslash}X|>{\centering\arraybackslash}X|>{\centering\arraybackslash}X|}
\hline
    对于 C++ 语言 & \texttt{\texttt{seq.cpp}} & \texttt{\texttt{demo2.cpp}} & \texttt{\texttt{demo3.cpp}} & \texttt{\texttt{demo4.cpp}} \\
\hline
\end{tabularx}
\end{center}

{\noindent\hspace*{1.5em}\rmfamily 编译选项}\par
\begin{center}
\begin{tabularx}{\textwidth}{|l|>{\centering\arraybackslash}X|>{\centering\arraybackslash}X|>{\centering\arraybackslash}X|>{\centering\arraybackslash}X|}
\hline
    对于 C++ 语言 & \multicolumn{4}{>{\centering\arraybackslash}X|}{\texttt{-O2 -std=c++14 -static}} \\
\hline
\end{tabularx}
\end{center}

{\noindent\hspace*{1.5em}\stress{注意事项（请仔细阅读）}}\par
\begin{enumerate}
    \item 文件名（程序名和输入输出文件名）必须使用英文小写。
    \item main 函数的返回值类型必须是 int，程序正常结束时的返回值必须是 0。
    \item 提交的程序代码文件的放置位置请参考各省的具体要求。
    \item 因违反以上三点而出现的错误或问题，申诉时一律不予受理。
    \item 若无特殊说明，结果的比较方式为全文比较（过滤行末空格及文本回车）。
\end{enumerate}
\end{titlepage}

\setcounter{page}{2}

\newpage

\section{ 数列之美（\englishname{ seq }）}

\subsection[ 题目背景 ]{【 题目背景 】}

\begin{callout}
“示例背景：这是一个虚构的题目，用于展示本模板的排版效果。”
\end{callout}

在遥远的数学王国里，数列们遵循着简单的规律生活。


\subsection[ 题目描述 ]{【 题目描述 】}

给定一个长度为 $n$ 的序列 $a_1, a_2, \ldots, a_n$，定义其\stress{「美丽值」}为

$$B = \sum_{i=1}^{n} a_i \cdot \operatorname{lowbit}(a_i)$$

请你求出 $B \bmod 998244353$ 的值。这里的 \texttt{lowbit} 指二进制下最低位的 $1$ 所对应的值。

\stress{数据保证}：所有 $a_i$ 均为正整数。


\subsection[ 输入格式 ]{【 输入格式 】}

输入文件：\filename{seq.in}

第一行一个正整数 $n$。

第二行 $n$ 个正整数 $a_i$。


\subsection[ 输出格式 ]{【 输出格式 】}

输出文件：\filename{seq.out}

输出一行一个整数，表示答案。


\subsection[ 样例 1 输入 ]{【 样例 1 输入 】}

\begin{minted}[linenos]{text}
3
1 2 3
\end{minted}


\subsection[ 样例 1 输出 ]{【 样例 1 输出 】}

\begin{minted}[linenos]{text}
8
\end{minted}


\subsection[ 样例 2 输入 ]{【 样例 2 输入 】}

\begin{minted}[linenos]{text}
5
7 7 7 7 7
\end{minted}


\subsection[ 样例 2 输出 ]{【 样例 2 输出 】}

\begin{minted}[linenos]{text}
35
\end{minted}


\subsection[ 样例 3 输入 ]{【 样例 3 输入 】}

\begin{minted}[linenos]{text}
1
2
3
4
5
6
7
8
9
10
11
12
13
14
15
16
17
18
19
20
21
22
23
24
25
26
27
28
29
30
31
32
33
34
35
36
37
38
39
40
41
42
43
44
45
46
47
48
49
50
51
52
53
54
55
56
57
58
59
60
61
62
63
64
65
66
67
68
69
70
71
72
73
74
75
76
77
78
79
80
\end{minted}


\subsection[ 样例 3 输出 ]{【 样例 3 输出 】}

\begin{minted}[linenos]{text}
1
\end{minted}


\subsection[ 提示 ]{【 提示 】}

对于 $100\%$ 的数据，满足 $1 \le n \le 10^5$，$1 \le a_i \le 10^9$。


\begin{center}\begin{tblr}{
  colspec = {c|c|c|c},
  hlines,
  hline{1} = {2pt},
  hline{2} = {1.5pt},
  hline{Z} = {2pt},
}
  子任务 & 分值 & $n \le$ & 特殊性质 \\
  $1$ & $10$ & $10$ & 无 \\
  $2$ & $20$ & $10^3$ & $a_i \le 100$ \\
  $3$ & $30$ & $10^5$ & $a_i$ 均为 $2$ 的幂 \\
  $4$ & $40$ & $10^5$ & 无 \\
\end{tblr}\end{center}


\subsection[ 样例 1 解释 ]{【 样例 1 解释 】}

样例 $1$：$B = 1 \cdot 1 + 2 \cdot 2 + 3 \cdot 1 = 8$。



\subsection[ 示例代码 ]{【 示例代码 】}

本题提供一份数据生成器作为参考，选手可在其基础上编写完整程序：

\begin{minted}[linenos, breaklines]{cpp}
#include <bits/stdc++.h>
using namespace std;
struct RandomNumberGenerator {
    unsigned long long state, a = 6364136223846793005ULL, c = 1442695040888963407ULL;
    explicit RandomNumberGenerator(unsigned long long seed) : state(seed) {}
    unsigned long long next() { return state = state * a + c; }
};
int main() {
    ios_with_stdio(false); cin.tie(nullptr);
    int n, k, seed; cin >> n >> k >> seed;
    RandomNumberGenerator rnd(seed);
    for (int i = 1; i <= n; ++i) cout << (rnd.next() % 1000000000ULL + 1) << (i == n ? '\n' : ' ');
    return 0;
}
\end{minted}


\subsection[ 提示 ]{【 提示 】}

\stress{注意}：答案需要对 $998244353$ 取模。

本题的评测环境为 Linux x86\_64，选手需要注意以下事项：

\begin{itemize}
  \item 使用 \texttt{long long} 类型存储中间结果；
\begin{itemize}
  \item 乘法结果可能超过 $2^{31}$；
  \item 取模运算应使用常量模数而非字面量。
\end{itemize}
  \item 避免在循环内进行不必要的 I/O 操作。
\end{itemize}



\newpage

\section{ 简单求和（\englishname{ demo2 }）}

\subsection[ 题目描述 ]{【 题目描述 】}

给定两个整数 $a$ 和 $b$，请输出它们的和 $a + b$。

这是一个用于演示最简题面结构的虚构题目。


\subsection[ 输入格式 ]{【 输入格式 】}

从标准输入中读入数据。

一行两个整数 $a, b$，用空格隔开。


\subsection[ 输出格式 ]{【 输出格式 】}

输出到标准输出中。

一行一个整数，表示 $a + b$。


\subsection[ 样例 1 输入 ]{【 样例 1 输入 】}

\begin{minted}[linenos]{text}
3 4
\end{minted}


\subsection[ 样例 1 输出 ]{【 样例 1 输出 】}

\begin{minted}[linenos]{text}
7
\end{minted}


\subsection[ 样例 2 输入 ]{【 样例 2 输入 】}

\begin{minted}[linenos]{text}
-5 10
\end{minted}


\subsection[ 样例 2 输出 ]{【 样例 2 输出 】}

\begin{minted}[linenos]{text}
5
\end{minted}



\newpage

\section{ 棋盘平方（\englishname{ demo3 }）}

\subsection[ 题目描述 ]{【 题目描述 】}

小 $\text{demo3}$ 有一个 $n \times m$ 的棋盘，每个格子上有一个正整数。

定义一次操作为：选择一行，将这一行所有格子上的数变为其平方。求经过若干次操作后，棋盘上所有数的最大值。


\subsection[ 输入格式 ]{【 输入格式 】}

从标准输入中读入数据。

第一行两个正整数 $n, m$。

接下来 $n$ 行，每行 $m$ 个正整数，表示棋盘上的数。


\subsection[ 输出格式 ]{【 输出格式 】}

输出到标准输出中。

输出一行一个整数，表示所有数的最大值。


\subsection[ 样例 1 输入 ]{【 样例 1 输入 】}

\begin{minted}[linenos]{text}
2 2
1 2
3 4
\end{minted}


\subsection[ 样例 1 输出 ]{【 样例 1 输出 】}

\begin{minted}[linenos]{text}
4
\end{minted}


\subsection[ 提示 ]{【 提示 】}

对于 $100\%$ 的数据，满足 $1 \le n, m \le 10^3$，格子上的数不超过 $10^4$。



\newpage

\section{ 回文判断（\englishname{ demo4 }）}

\subsection[ 题目描述 ]{【 题目描述 】}

给定一个字符串 $s$，请判断它是否是回文串。

这是一个用于演示极简题面的虚构题目。


\subsection[ 输入格式 ]{【 输入格式 】}

从标准输入中读入数据。

一行一个字符串 $s$。


\subsection[ 输出格式 ]{【 输出格式 】}

输出到标准输出中。

若 $s$ 是回文串输出 \texttt{Yes}，否则输出 \texttt{No}。


\subsection[ 样例 1 输入 ]{【 样例 1 输入 】}

\begin{minted}[linenos]{text}
abcba
\end{minted}


\subsection[ 样例 1 输出 ]{【 样例 1 输出 】}

\begin{minted}[linenos]{text}
Yes
\end{minted}


\subsection[ 样例 2 输入 ]{【 样例 2 输入 】}

\begin{minted}[linenos]{text}
hello
\end{minted}


\subsection[ 样例 2 输出 ]{【 样例 2 输出 】}

\begin{minted}[linenos]{text}
No
\end{minted}


\end{document}